\documentclass[authoryear,preprint,review,12pt]{elsarticle}


\textheight24cm
\textwidth15.5cm              
\oddsidemargin0.25cm
\topmargin-1.8cm
\usepackage[utf8]{inputenc}
\usepackage{booktabs} 
\usepackage{tabularx} 
\usepackage{ltablex} 
\usepackage{geometry} 
\usepackage{caption}
\usepackage{amsmath}
\usepackage{tcolorbox}
\usepackage{soul}


\usepackage{caption}
% Adjust the page margins to fit the table size
\geometry{
  left=0.75in,
  right=0.75in,
  top=1in,
  bottom=1in
}


\usepackage[T1]{fontenc}
\usepackage{lmodern}
\usepackage{pifont}
\usepackage{placeins}
\usepackage{color}
\newcommand{\red}[1]{\textcolor{red}{#1}}
\usepackage{graphicx}
\usepackage{pdfpages}
\usepackage{lineno}  
\usepackage{subcaption} 
\usepackage{amsmath} 
\usepackage{amssymb}
\usepackage[utf8]{inputenc}
\usepackage{float}
\usepackage{arydshln}
\usepackage{url}
\def\UrlBreaks{\do\/\do-}
\usepackage{multirow}
\usepackage{comment}
\usepackage{threeparttable}
\usepackage{bm}
\usepackage{amsmath}
\usepackage{booktabs} 
\usepackage{tabularx}
\usepackage{geometry} 

%\usepackage[x11names]{xcolor}
\usepackage[
  colorlinks,
]{hyperref}
\definecolor{DodgerBlue4}{rgb}{0.0627, 0.3059, 0.5451}
\AtBeginDocument{\hypersetup{citecolor=DodgerBlue4}}

%%%%%%%%%%%%%%%%%%%%%%%%%%%%%%%%%%%%%%%%%%%%%%%%%%%
\journal{International Journal of Geographical Information Science}

\begin{document}
\begin{frontmatter}
\title{Measuring causal strengths by geographical cross mapping cardinality}

\author[a,b,c]{Wenbo Lv}
\ead{wenbolyu@hkust-gz.edu.cn}

\author[d]{Shaoqing Dai}
\ead{shaoqing.dai@whu.edu.cn}

\author[e]{Yongze Song\corref{cor1}}
\cortext[cor1]{Corresponding author: Yongze Song, yongze.song@curtin.edu.au}

\author[a]{Wufan Zhao}
\ead{wufanzhao@hkust-gz.edu.cn}

\author[b]{Wen Yi}
\ead{wen.yi@polyu.edu.hk}

\author[a]{Yumiao Xiao}
\ead{yxiao224@connect.hkust-gz.edu.cn}

\author[f,g]{Nan Jia}
\ead{jianan2@msu.edu}

\address[a]{Urban Governance and Design Thrust, Society Hub, The Hong Kong University of Science and Technology(Guangzhou), Guangzhou, China}
\address[b]{Department of Building and Real Estate, The Hong Kong Polytechnic University, Hong Kong, China}
\address[c]{The Hong Kong Polytechnic University, Shenzhen Research Institute, Shenzhen, China}
\address[d]{School of Resource and Environmental Sciences, Wuhan University, Wuhan, China}
\address[e]{School of Design and the Built Environment, Curtin University, Perth, Australia}
\address[f]{Department of Fisheries and Wildlife, Michigan State University, East Lansing, USA}
\address[g]{School of Design, Shanghai Jiao Tong Univeristy, Shanghai, China}


%%%%%%%%%%%%%%%%%%%%%%%%%%%%%%%%%%%%%%%%%%%%%%%%%%%
\begin{abstract}


Spatial cross-sectional data encapsulate rich information on spatial processes, serving as a critical foundation for examining causation between variables. Detecting and quantifying these causation is essential for understanding the mechanisms of complex natural and human phenomena. However, measuring causal strengths from spatial cross-sectional data remains challenging, as existing methods lack a rigorous causal framework to quantify the causation. Here, we propose a geographical cross mapping cardinality (GCMC) model, which measures the causal strengths using the intersectional cardinality of neighbors identified through mutual cross mapping in the reconstructed state space, while integrating the DeLong placement method to test the statistical significance and confidence intervals of the causal strength. We validate the applicability of GCMC by applying it to three typical spatial cross-sectional datasets with clearly existing causations and evaluating its sensitivity to noise in observed data. The results indicate that GCMC can effectively measure causal strengths in weakly, moderately, and strongly coupled systems while also revealing the dominant(primary) causal direction. Moreover, GCMC exhibits robustness to noise in observational data. GCMC is a new method within empirical dynamic modeling frameworks for spatial cross-sectional data and can be used alongside existing methods to more robustly infer causation.


\end{abstract}

\begin{keyword}
Causal Strength \sep Spatial Cross-Sectional Data \sep Empirical Dynamic Modeling \sep Intersectional Cardinality \sep Geographical Cross Mapping Cardinality
\end{keyword}

\end{frontmatter}

%%%%%%%%%%%%%%%%%%%%%%%%%%%%%%%%%%%%%%%%%%%%%%%%%%%
\linenumbers


 \section{Introduction}

Understanding causation in the Earth system is fundamental for addressing scientific and societal challenges, yet direct interventions are often impractical or ethically problematic \citep{gao2022tscess,runge2023tsc,gao2023gccm}. Traditionally, Earth system sciences have relied on physics-based simulations to infer causal mechanisms, offering a controlled framework for experimentation. However, these models come with high computational costs and rely on strong assumptions about underlying processes, limiting their applicability in some contexts. Meanwhile, the increasing availability of large-scale observational datasets, including time series and spatial cross-sectional data, has opened new avenues for data-driven causal inference in Earth sciences \citep{runge2019esstsc,runge2023tsc}. A variety of methods based on observational data have been developed, including prediction-based approaches that test temporal precedence and predictive power \citep{granger1969grangertest,geweke1982grangertest2}, entropy-based techniques that capture nonlinear dependencies \citep{schreiber2000te1,vicente2010te2,sun2014ce1}, structural causal modeling frameworks that encode domain knowledge into graphical models \citep{pearl2000scm2,pearl2009scm3,pearl2018scm4}, potential outcome frameworks that enable counterfactual reasoning in quasi-experimental designs \citep{rubin1974pom1,rubin1986pom2,rubin2005pom3}, and empirical dynamic modeling approaches that detect causation through delay-coordinate embedding and attractor reconstruction in nonlinear dynamic systems \citep{sugihara2012ccm,ma2014cms,clark2015multispatialccm}.

Causal inference in Earth system science still faces several challenges, despite numerous algorithmic models for causal inference based on observational data have been developed. Most of these causation models are designed for time series data, such as Peter Clark algorithm with momentary conditional independence test (PCMCI) \citep{runge2019pcmci}, convergent cross mapping (CCM) \citep{sugihara2012ccm}, and partial cross mapping (PCM) \citep{leng2020pcm}, and there remains a lack of causation models specifically tailored for spatial cross-sectional data \citep{gao2022ciss,gao2023gccm}. Compared to time series data, spatial cross-sectional data also capture rich information about Earth system processes. \cite{gao2022tscess} suggested that when temporal variations are not prominent, spatial variations can be used to infer causation in Earth systems, for the need to fully utilize spatial cross-sectional data for causal inference. Despite this potential, a significant gap remains: most existing causal inference methods are not explicitly designed to handle the unique characteristics of spatial cross-sectional data. Therefore, it is essential to evaluate whether current spatial analytical models can reliably infer causation, or whether new frameworks grounded in formal causal principles are needed.

The existing models commonly used for causal inference in spatial cross-sectional data, such as pearson correlation coefficient (PCC) \citep{pearson1895pcc}, spatial autoregression \citep{anselin2013sar}, geographically weighted regression \citep{fotheringham2003gwr,fotheringham2017mgwr}, and geographical detector (GD) \citep{wang2010gd1,lyu2025gdverse}, are not strictly designed based on formal causal principles. By causal principles, we refer to frameworks that aim to distinguish causation from correlation through explicit assumptions or structures, including: (i) the counterfactual framework (as in the Rubin Causal Model), (ii) structural causal models using directed acyclic graphs and do-calculus (as in Pearl’s framework), and (iii) predictive causality based on temporal or dynamical predictability, such as Granger causality and its nonlinear extensions. Although the aforementioned spatial models are widely used to interpret relationships with causal implications, they do not inherently identify causal directions or account for confounding factors, and thus may yield spurious conclusions when applied to observational spatial data \citep{gao2022ciss}. In addition, some spatialized causal inference methods incorporate formal causal frameworks, such as spatial difference-in-differences \citep{delgado2015sdid} and spatial regression discontinuity designs \citep{kolak2019srdd}. However, these methods still face fundamental limitations when applied to spatial cross-sectional data. For instance, the presence of spatial spillover effects often violates the stable unit treatment value assumption (SUTVA), and spatial heterogeneity complicates the identification of localized treatment effects \citep{akbari2021spatialcausality,gao2023gccm}, making it difficult to extract valid causal information from spatial observations alone.

The Earth system is a complex dynamical system with deeply interconnected components \citep{runge2019esstsc,runge2023tsc}, where classical assumptions of static or separable units often break down. In this context, causal inference methods based on predictive capability, such as empirical dynamic modeling (EDM), provide a promising alternative. EDM infers causality by examining whether the historical states of one variable can be used to reconstruct and predict the states of another, reflecting information flow in dynamic systems rather than relying on correlation alone \citep{sugihara2012ccm,leng2020pcm,tao2023cmc}. This dynamic perspective aligns with the notion of Granger causality, but extends it to nonlinear and complex systems.

Recent efforts have extended EDM to spatial domains. Notably, the Geographical Convergent Cross Mapping (GCCM) model adapts the original CCM method \citep{sugihara2012ccm} to spatial cross-sectional data by leveraging state-space reconstruction techniques \citep{lorenz1969ssr,Takens1981embedding,Mane1981embedding}. This pioneering approach enables the inference of causal links from spatial snapshots, addressing longstanding challenges in spatial causal inference that traditional temporal EDM methods cannot resolve \citep{gao2022tscess,gao2023gccm}.

Building on these developments, this study proposes a novel model, Geographical Cross Mapping Cardinality (GCMC), for quantifying causal strengths between variables in spatial cross-sectional data. Inspired by Cross Mapping Cardinality (CMC), which uses the area under the intersectional cardinality curve (AUC) to measure causal strength \citep{tao2023cmc}, GCMC constructs spatially reconstructed state spaces and evaluates the intersectional cardinality (IC) curves between variable manifolds. To statistically estimate the causal strength and its confidence interval, we apply the DeLong placement method \citep{delong1988,sun2014delongfi} to the IC curve. We demonstrate the effectiveness of GCMC using a benchmark system consisting of three variables in canonical causal interaction structures \citep{sanchez2024surd}: mediator ($A \rightarrow B \rightarrow C$), collider ($A \rightarrow B \leftarrow C$), and confounder ($A \leftarrow B \rightarrow C$); three real-world case studies covering distinct coupling states; and further assess its robustness against observational noise using a synthetic dataset. Comparative analysis with PCC, GD, and GCCM confirms the practical advantages of GCMC.

The remainder of this manuscript is structured as follows: Section~\ref{sec2} outlines the theoretical foundations and methodology of the proposed GCMC method, including its rationale and validation design. Section~\ref{sec3} reports the experimental results from synthetic and real-world datasets. Section~\ref{sec4} discusses key findings, and Section~\ref{sec5} concludes this study.

\section{Material and methods}\label{sec2} 

\subsection{The geographical cross mapping cardinality method}

Generalized embedding theory \citep{Takens1981embedding,Mane1981embedding} posits that the dynamics of a system can be reconstructed from a single observable through delay-coordinate embedding. Building upon this foundation, we extend the theory to the spatial domain by using spatial neighbors as analogs to temporal delays, thereby constructing shadow manifolds that capture the intrinsic structure of the system. Let $x$ and $y$ denote two spatially distributed (spatial cross-sectional) variables with $n$ observations each, the reconstructed state space $E_x$ of variable $x$ is given by:
\begin{equation}\label{eq-embedding}
E_x = \begin{bmatrix}
h_{s(1)}(x_1) & h_{s(1+\tau)}(x_1) & \cdots & h_{s(1+(E-1)\tau)}(x_1) \\
h_{s(1)}(x_2) & h_{s(1+\tau)}(x_2) & \cdots & h_{s(1+(E-1)\tau)}(x_2) \\
\vdots & \vdots & \ddots & \vdots \\
h_{s(1)}(x_n) & h_{s(1+\tau)}(x_n) & \cdots & h_{s(1+(E-1)\tau)}(x_n)
\end{bmatrix}
\end{equation}
\noindent where $E$ denotes the embedding dimension; $\tau$ is the spatial lag step; $s(k)$ identifies the $k$-th order spatial neighbors of each unit; and $h(\cdot)$ aggregates attribute values from neighboring units into scalar using the mean. 

For instance, as illustrated in Figure~\ref{fig_ssr}, when the embedding dimension \(E = 3\) and the spatial lag step \(\tau = 1\), the embedding vector for a focal unit \(i\) is computed as 
\(h_{s(1)}(x_i) = \frac{1145 + 749 + 749 + 972 + 625}{5} = 848\),  
\(h_{s(2)}(x_i) = \frac{692 + 688 + 479 + 387 + 430 + 885}{6} = 593.5\), and  
\(h_{s(3)}(x_i) = \frac{734 + 698 + 859 + 323 + 225 + 357}{6} \approx 532.67\),  
where the denominators correspond to the number of units included in each spatial lag calculation and the numerators represent the sum of the attribute values of the neighbors at the corresponding lag order.

\begin{figure}[htb!]
\centering
\includegraphics[width=1\textwidth]{./figure/figure1.jpg}
\caption{Demonstration of reconstructing state space for spatial cross-sectional data.}
\label{fig_ssr}
\end{figure}

If $x$ and $y$ are governed by a common underlying dynamical mechanism, their reconstructed state spaces (or manifolds) will exhibit structural overlap. This overlap enables mutual cross mapping, forming the basis for causal inference. This principle is rooted in the theory of dynamical systems and generalized embedding theory \citep{Takens1981embedding,Mane1981embedding,sugihara2012ccm}, which asserts that if two variables originate from the same system, their reconstructed phase-space trajectories will reflect shared dynamics. In such cases, information about one variable can be reliably inferred from the other, indicating a potential causal relationship. Once the state-space manifolds $E_x$ and $E_y$ are reconstructed, GCMC evaluates the consistency of mutual cross mapping between them to infer causality. For a given number of nearest neighbors $k$, let $E_y^{nn}(i,k)$ denote the set of the $k$ nearest neighbors of the $i$-th spatial unit in $E_y$, and $E_x^{nn}(i,k)$ the corresponding set in $E_x$. By projecting the neighbors from $E_x$ to $E_y$, we obtain $E_y^{nn}(E_x^{nn}(i,k),k)$—that is, the set of $k$ nearest neighbors in $E_y$ corresponding to the $k$ nearest neighbors of unit $i$ in $E_x$. This projection reflects the idea of using the local structure in $x$ to reconstruct the neighborhood of $y$, a hallmark of causality in nonlinear dynamical systems.

To quantify the consistency between these neighborhoods, we define the intersection cardinality for the $i$-th unit as:
\begin{equation}\label{eq-ic_ik}
IC_{i,k} = \left| E_y^{nn}(i,k) \cap E_y^{nn}(E_x^{nn}(i,k),k) \right|
\end{equation}
\noindent This value measures the number of shared neighbors between the directly observed neighborhood in $E_y$ and the neighborhood inferred via cross mapping from $E_x$. A key insight from CMC \citep{tao2023cmc} is that a higher IC indicates stronger alignment between the true and predicted neighborhood structures, suggesting a more robust causal link from $y$ to $x$.

Averaging across all spatial units yields the global IC value for a given neighborhood size $k$:
\begin{equation}\label{eq-ic_k}
IC_k = \frac{1}{n} \sum_{i=1}^n IC_{i,k}
\end{equation}

By varying $k$ from $1$ to $n$, we obtain the complete IC curve:
\begin{equation}\label{eq-ic_curve}
IC_{\text{curve}} = \{ IC_k \mid k = 1, 2, \dots, n \}
\end{equation}

According to \cite{tao2023cmc}, the shape of the IC curve reveals critical information about the presence and strength of causality. Under the null hypothesis $H_0$(Figure~\ref{fig_ic}b,which assumes no causal influence from $y$ to $x$), the IC curve increases linearly with $k$, reflecting the expected number of intersecting neighbors under random mapping. In contrast, under the alternative hypothesis $H_1$(Figure~\ref{fig_ic}a), where a causal link exists, the IC curve exhibits a concave-down shape, indicating improved neighborhood consistency with increasing $k$. This IC concavity serves as a dynamical signature of causality. Formally, to test for causality from $y$ to $x$, we specify the hypotheses as:
\begin{equation}\label{eq-hypothesis}
H_0: y \not\Rightarrow x \quad \text{vs.} \quad H_1: y \Rightarrow x
\end{equation}

The IC curve under $H_1$ is obtained through equations~\ref{eq-ic_ik}–\ref{eq-ic_curve}, whereas the IC curve under $H_0$ is approximated by a baseline linear function, $IC_k = 0.5k$, representing the expected overlap under random association \citep{tao2023cmc}. The causal strength (CS) from $y$ to $x$ is then defined as the difference in the area under the IC curve (AUC) between the alternative and null hypotheses(Figure~\ref{fig_ic}c):
\begin{equation}\label{eq-cs}
CS_{y \rightarrow x} = AUC(H_1 \mid H_0) = AUC(H_1) - AUC(H_0)
\end{equation}

\begin{figure}[htb!]
\centering
\includegraphics[width=1\textwidth]{./figure/figure2.jpg}
\caption{Schematic diagram of intersectional cardinality (IC) in GCMC. \textbf{a} The intersectional cardinality corresponding to the \(H_1\) model, which indicates the existence of a causation when neighbors map to neighbors. \textbf{b} The intersectional cardinality corresponding to the \(H_0\) model, where neighbors map to random points, shows no significant causation. \textbf{c} The intersectional cardinality curve, formed by the intersectional cardinality values corresponding to \(H_1\) and \(H_0\) for different numbers of nearest neighbors \(k\)}
\label{fig_ic}
\end{figure}

In GCMC, the causal strength (AUC value) is estimated using the DeLong placement method, which is a non-parametric approach for calculating the AUC and its variance \citep{delong1988,sun2014delongfi}. Given a set of positive cases \(\{X_i\}_{i=1}^{m}\) (IC curves corresponding to \(H_1\)) and negative controls \(\{Y_j\}_{j=1}^{n}\) (IC curves corresponding to \(H_0\)), the AUC is computed as:

\begin{equation}
\hat{\theta} = \frac{1}{mn} \sum_{i=1}^{m} \sum_{j=1}^{n} \mathbb{I}(X_i > Y_j) + \frac{1}{2} \mathbb{I}(X_i = Y_j),
\end{equation}

\noindent where \(\hat{\theta}\) represents the estimated AUC value, and \(\mathbb{I}(\cdot)\) is an indicator function that equals 1 if the condition holds and 0 otherwise.  

To assess the statistical significance and confidence interval of the estimated AUC, the DeLong placement method provides a variance estimate based on the placement values \citep{delong1988,sun2014delongfi}:

\begin{equation}
\hat{\sigma}^2 = \frac{1}{m} \sum_{i=1}^{m} \left( X_i - \hat{\theta} \right)^2 + \frac{1}{n} \sum_{j=1}^{n} \left( Y_j - (1 - \hat{\theta}) \right)^2.
\end{equation}

Using this variance, a confidence interval for AUC can be constructed based on the normal approximation:

\begin{equation}
\text{CI} = \left[ \hat{\theta} - z_{\alpha/2} \hat{\sigma}, \hat{\theta} + z_{\alpha/2} \hat{\sigma} \right],
\end{equation}

\noindent where \( z_{\alpha/2} \) is the critical value of the standard normal distribution for a given confidence level (e.g., 1.96 for 95\% confidence), and \(\hat{\sigma} \) is the estimated standard error of \(\hat{\theta}\), which quantifies the uncertainty in the estimated AUC value.

To evaluate the statistical significance of \(\hat{\theta}\) (the AUC value), a hypothesis test is conducted to determine whether the AUC significantly deviates from 0.5 under the null hypothesis \(\theta = 0.5\), where the test statistic is given by

\begin{equation}
Z = \frac{\hat{\theta} - 0.5}{\hat{\sigma}}.
\end{equation}

The corresponding two-tailed p-value is obtained using the standard normal cumulative distribution function:

\begin{equation}
p = 2 \Phi(-|Z|),
\end{equation}

\noindent where \(\Phi(\cdot)\) is the cumulative distribution function of the standard normal distribution.

The causal strength measured by GCMC is a metric rooted in dynamical systems theory, quantifying the difference between intersection cardinality curves under the null and alternative hypotheses of mutual cross mapping. To reliably infer a causal link, GCMC requires that the estimated causal strength in a given direction be both statistically significant and sufficiently large in magnitude. In bivariate systems, if neither direction yields significant causal strength, it suggests the absence of a causal relationship between the variables. If one or both directions show significance, prior domain knowledge is essential to determine whether the measured causal strength represents a genuine causal influence or arises from strong coupling or synchronization (e.g., enslaved effects), which may lead to false positives. In multivariate settings, GCMC evaluates all possible directional pairs, retaining only those that meet both the statistical significance criterion and exceed a predefined threshold (e.g., 0.1). This dual criterion—statistical significance and threshold filtering—ensures that the identified causal directions are robust and interpretable. As a result, GCMC offers a principled and reliable approach to causal inference, with enhanced resistance to false positives compared to existing methods.

\subsection{Validation strategy and experimental design}

\subsubsection{Synthetic benchmark via spatial logistic map}

To evaluate the capability of the proposed GCMC method in measuring causal strengths from spatial cross-sectional data, we design a synthetic benchmark system inspired by causal simulation practices in nonlinear time series. In particular, we extend the classic logistic map into a spatial domain to construct a \textit{spatial logistic map} process \citep{willeboordse2003slm}, enabling the controlled simulation of structured causal interactions over spatial units. This spatial benchmark is used to generate three prototypical causal structures among three variables \(A\), \(B\), and \(C\): \textbf{mediator} $A \rightarrow B \rightarrow C$, \textbf{collider} $A \rightarrow B \leftarrow C$, and \textbf{confounder} $A \leftarrow B \rightarrow C$ (Figure~\ref{fig_benchmark}a), which are commonly considered in causal inference literature \citep{sanchez2024surd}.

\begin{figure}[htb!]
\centering
\includegraphics[width=1\textwidth]{./figure/figure3.jpg}
\caption{Benchmark evaluation of GCMC on canonical causal structures.
\textbf{(a)} Three simulated causal scenarios representing fundamental three-variable causal motifs. \textbf{(b)} Performance comparison of GCMC and GCCM on each scenario using estimated causal strengths (Values near edges). Edge types indicate causal inference correctness: true positives (black solid), false positives (red solid), true negatives (gray dashed), and false negatives (red dashed). Non-significant values are marked with $^\#$ and orange-highlighted thresholds indicate the optimal decision boundary.}
\label{fig_benchmark}
\end{figure}

Each synthetic spatial variable is initialized using a stochastic spatial field generated over a regular 2D grid. Specifically, we construct a distance matrix across all spatial units and apply an exponential covariance function \(\text{Cov}(d) = \sigma^2 \exp(-d / \rho)\), where \(d\) is the Euclidean distance between spatial locations, \(\sigma^2 = 1\) is the marginal variance, and \(\rho = 1.5\) is the range parameter controlling spatial autocorrelation. Spatially correlated fields are then drawn from a multivariate normal distribution parameterized by this covariance matrix. To ensure non-negativity and comparability across simulations, the resulting fields are truncated at zero to remove negative values and rescaled to the unit interval \([0,1]\). Each of the three variables \(A\), \(B\), and \(C\) is initialized independently through this procedure, producing realistic spatial patterns with inherent autocorrelation.

We then evolve the system using the spatial logistic map for 15 time steps. The update rule for any variable \(V \in \{A, B, C\}\) at location \(\mathbf{s}_i\) and time \(t+1\) is given by:
\begin{equation}\label{slm}
V_{t+1}(\mathbf{s}_i) = 1 - \alpha_V V_t(\mathbf{s}_i)\left(\frac{1}{k}\sum_{j\in \mathcal{N}(i)} V_t(\mathbf{s}_j) - \sum_{U\in \mathcal{C}(V)}\beta_{UV} U_t(\mathbf{s}_i)\right)
\end{equation}
\noindent where \(k\) is the neighborhood size; \(\mathcal{N}(i)\) denotes the spatial neighborhood of unit \(i\); \(\alpha_V\) controls the strength of nonlinear spatial self-interaction via a logistic growth term; \(\beta_{UV}\) denotes the causal influence from variable \(U\) to \(V\); and \(\mathcal{C}(V)\) denotes the set of causal parent variables of \(V\).

To generate spatial cross-sectional data for testing, we discard the transient initial state at \(t=0\) and compute the average value of each variable over all spatial units at each subsequent time step. The final test data are obtained by averaging these spatial unit values across the simulation steps after the initial transient period, thereby capturing the stabilized spatial patterns. The parameters for the three causal scenarios are set as follows: for the \textbf{mediator} case, we set \(\beta_{AB}=1\), \(\beta_{BC}=1\), and all other \(\beta\) coefficients are set to $0$; for the \textbf{collider} case, \(\beta_{AB}=1\), \(\beta_{CB}=1\), with all remaining \(\beta\) coefficients equal to $0$; and for the \textbf{confounder} case, \(\beta_{BA}=1\), \(\beta_{BC}=1\), while all other \(\beta\) coefficients are $0$. In all scenarios, the spatial logistic parameters are fixed at \(\alpha_A = \alpha_B = \alpha_C = 0.2\), and the spatial neighborhood size is set to \(k=4\).

\subsubsection{Case studies in earth and social systems}

The developed GCMC model is validated through three real-world case studies covering diverse socio-economic and environmental contexts to comprehensively assess its practical applicability and comparative performance under realistic spatial conditions. The first case employs the Columbus Neighborhood Data from the Columbus, Ohio metropolitan area \citep{r-spData}, examining causal relationships among housing value (\textit{hoval}), household income (\textit{inc}), and crime rates (\textit{crime}, measured by residential burglaries and vehicle thefts per thousand households). The second case analyzes county-level data across mainland China \citep{gao2023gccm}, investigating potential causal links among population density (\textit{popd}), elevation (\textit{elev}), and temperature (\textit{tem}). The third case focuses on farmland productivity in China, using the same data source \citep{gao2023gccm} to explore causal interactions among net primary productivity (\textit{npp}), precipitation (\textit{pre}), and temperature (\textit{tem}). Across all cases, causal inference results generated by the GCMC model are benchmarked against PCC, GD, and GCCM to evaluate empirical validity and robustness. For the GD model, the continuous variables are discretized into five categories using the natural breaks classification method. For both GCMC and GCCM, the optimal embedding dimension \(E\) is determined using the false nearest neighbors (FNN) method \citep{kennel1992fnn}. During GCMC execution, the number of neighbors is selected from a candidate set based on the value of neighbors that yields the highest causal strength. In contrast, GCCM uses \(E + 2\) nearest neighbors for cross mapping. Comprehensive visualizations accompany the analyses to support thorough interpretation of the results.

\subsubsection{Noise sensitivity analysis}

To evaluate the robustness of the GCMC model against observational noise, a controlled noise sensitivity experiment is conducted. We focus on the variable pair \textit{popd} and \textit{elev} from the population density case study. Noise is introduced to these variables following the formulations:

\begin{equation}\label{eq-sensitivity}
\begin{aligned}
S_{popd} &= popd + a \cdot \eta \cdot popd, \\
S_{elev} &= elev + a \cdot \eta \cdot elev,
\end{aligned}
\end{equation}

where \(popd\) and \(elev\) represent the original data vectors, \(a\) denotes the noise ratio varying from 0 to 1 in increments of 0.1, and \(\eta\) is a standard normal random vector matching the length of the variables. The resulting \(S_{popd}\) and \(S_{elev}\) correspond to the noisy versions of the original variables. The noise sensitivity of GCMC is assessed under three scenarios: noise added to \textit{elev} only, noise added to \textit{popd} only, and noise simultaneously added to both variables.

\section{Results}\label{sec3} 

\subsection{Performance on benchmark systems}\label{benchmark}

Figure~\ref{fig_benchmark} illustrates the performance of GCMC and GCCM on three canonical three-variable causal configurations: mediator, collider, and confounder structures. In each scenario (Figure~\ref{fig_benchmark}a), variables \(A\), \(B\), and \(C\) interact according to a distinct causal topology, and Figure~\ref{fig_benchmark}b presents the causal strengths inferred by both methods, with arrows indicating detected directions and numeric labels denoting strength values. In the mediator case ($A \rightarrow B \rightarrow C$), both methods successfully identify the causal link from \(A\) to \(B\), and both detect a causal link from \(A\) to \(C\), which is theoretically justified due to the transitive effect of \(A\) on \(C\) mediated through \(B\) \citep{leng2020pcm}; however, only GCMC correctly detects the causal link from \(B\) to \(C\), while GCCM misses this link. In the collider case ($A \rightarrow B \leftarrow C$), GCMC identifies only the two true causes of \(B\), whereas GCCM erroneously infers a link from \(A\) to \(C\), likely reflecting collider bias. In the confounder configuration ($A \leftarrow B \rightarrow C$), both methods detect the true outgoing links from \(B\), but only GCCM incorrectly infers a nonexistent link from \(A\) to \(C\), which GCMC avoids. Overall, GCMC achieves higher causal inference accuracy than GCCM across all tested structures, demonstrating its effectiveness in reducing both false positives and false negatives under nonlinear spatial causal configurations.

\subsection{Case of residential crime study}\label{case1}

\begin{figure}[htbp]
\centering
\includegraphics[width=1\textwidth]{./figure/figure4.jpg}
\caption{Causation among residential crime, house value, and household income (non-significant values are marked with $^\#$). \textbf{a-c} Spatial distributions of the three variables. \textbf{d} Correlations obtained by PCC, where high correlations indicate a moderately to strongly coupled scenario. \textbf{e} Stratified power computed by GD. \textbf{f} Causal strengths measured by GCCM. \textbf{g} Causal strengths measured by GCMC. \textbf{h} Conceptual diagram summarizing the inferred causal pathways based on GCMC estimates, where false positive links are excluded based on a causal strength threshold of $0.1$. Both GCMC and GCCM were configured with an embedding dimension \(E = 7\); the number of nearest neighbors used was \(k = 25\) for GCMC and \(k = 9\) for GCCM.}
\label{fig_case1}
\end{figure}

This case study investigates the causal relationships among residential crime (crime), housing value (hoval), and household income (inc) using spatial data from Columbus, Ohio. The spatial distributions of the three variables (Figure~\ref{fig_case1}a--c) reveal strong urban--suburban contrasts. The correlation analysis (Figure~\ref{fig_case1}d) indicates moderate to strong negative associations between crime and both housing value and income, but does not provide directional insights into causation. GD (Figure~\ref{fig_case1}e) captures moderate stratified power between crime and hoval, as well as high stratified power between crime and inc, yet fails to distinguish between causal drivers and response variables. In contrast, GCMC (Figure~\ref{fig_case1}g) successfully identifies two actual causal links: from crime to hoval and from inc to hoval. Other links, though statistically significant, are filtered out as false positives based on the causal strength threshold of 0.1. Compared to GCMC, GCCM (Figure~\ref{fig_case1}f) additionally detects a spurious bidirectional causal relationship between inc and crime, despite the absence of theoretical or empirical evidence supporting such causation. These findings highlight the superior capability of GCMC in suppressing false positive causal links and improving the robustness and interpretability of causal inference in spatial observational data.

\subsection{Case of population density study}\label{case2}

\begin{figure}[htb]
\centering
\includegraphics[width=0.9\textwidth]{./figure/figure5.jpg}
\caption{Causation among elevation, temperature and population density (non-significant values are marked with $^\#$). \textbf{a} Correlations derived from PCC, which indicate that the dataset in this case represents a weak to moderate coupling state. \textbf{b} The stratified power quantified with GD. \textbf{c} Causal strengths measured by GCCM. \textbf{d} Causal strengths measured by GCMC. \textbf{e} Diagram of causal links constructed from GCMC results, with false positives filtered using a causal strength threshold of $0.1$. Both GCMC and GCCM were configured with an embedding dimension \(E = 10\); the number of nearest neighbors used was \(k = 200\) for GCMC and \(k = 12\) for GCCM.}
\label{fig_case2}
\end{figure}

As illustrated in Figure~\ref{fig_case2}, this case reveals distinct capabilities and limitations of existing methods in identifying causal relationships among elevation, temperature, and population density. While PCC (Figure~\ref{fig_case2}a) captures only linear correlations without indicating causality, GD (Figure~\ref{fig_case2}b) quantifies spatial stratified heterogeneity and provides limited evidence of potential influence. GCCM (Figure~\ref{fig_case2}c) measures strong causal strengths but erroneously suggests that temperature has a greater causal influence on elevation than the other way around - an unrealistic reversal given the known topographic control on temperature. In contrast, GCMC (Figure~\ref{fig_case2}d) successfully identifies the correct unidirectional causal link from elevation to temperature, as the causal strength from temperature to elevation is not statistically significant, while the causal strength from elevation to temperature exceeds the predefined threshold of 0.1, aligning with physical understanding. The resulting causal diagram (Figure~\ref{fig_case2}e) clearly illustrates the causal structure, demonstrating the ability of GCMC to distinguish true causal pathways and avoid spurious inferences, thus offering a more reliable and interpretable framework for causal discovery in spatial cross-sectional data.

\subsection{Case of net primary productivity study}\label{case3}

\begin{figure}[htb!]
\centering
\includegraphics[width=0.9\textwidth]{./figure/figure6.jpg}
\caption{Causation among temperature, precipitation and net primary productivity (non-significant values are marked with $^\#$). \textbf{a} Correlations measured by PCC, which suggest that this case exhibits a high degree of coupling. \textbf{b} Stratified power computed by GD. \textbf{c} Causal strengths measured by GCCM. \textbf{d} Causal strengths measured by GCMC. \textbf{e} Causal links inferred by GCMC after filtering false positives using a causal strength threshold of $0.21$. Both GCMC and GCCM were configured with an embedding dimension \(E = 18\); the number of nearest neighbors used was \(k = 165\) for GCMC and \(k = 20\) for GCCM.}
\label{fig_case3}
\end{figure}

For the well-known Climate-NPP causation, where classical time series causal models failed to capture the causation \citep{gao2022tscess,gao2023gccm}, we also tested whether GCMC can effectively measure its causal strengths. The relevant results are presented in Figure~\ref{fig_case3} and the true causation of precipitation and temperature on NPP were detected by GCMC, with a causal strength of $0.2397$ from precipitation to NPP (pre  \(\rightarrow\) npp) and $0.4288$ from temperature to NPP (tem  \(\rightarrow\) npp). The causal strength from NPP to temperature (npp \(\rightarrow\) tem) did not pass the significance test ($p > 0.05$), indicating that NPP does not exert a causal influence on temperature. Instead, this suggests that the measured causal strength by GCMC is a result of the enslaved effect of the strong causal influence of temperature on NPP. Compared to GCMC, GCCM additionally identified a spurious causal link from NPP to precipitation (npp \(\rightarrow\) pre), which was not supported by GCMC (false links were filtered using a threshold of $0.21$), indicating its susceptibility to false positives under strong coupling. In this case, GCMC successfully identified the actual causation, whereas the GD model identifies stratified powers where NPP shows a stronger linkage to precipitation and temperature than the reverse. This suggests that the stratified powers derived from the GD model do not provide a reliable indication of the actual causal direction. The results of this case indicate that under strong coupling, GCMC can effectively measure the actual causation, while also providing support for diagnosing enslaved effect. In this case, the data is organized in a spatial grid format, whereas the previous two cases use spatial lattice data, thereby demonstrating that GCMC is applicable not only to spatial lattice data but also to spatial grid data.

\subsection{Sensitivity to observational noise}\label{sensitivity}

\begin{figure}[htb]
\centering
\includegraphics[width=1\textwidth]{./figure/figure7.jpg}
\caption{Results of GCMC applied to synthetic datasets with varying noise levels. \textbf{a} Causal strengths when noise is added only to \textit{elev}. \textbf{b} Causal strengths when noise is added only to \textit{popd}. \textbf{c} Causal strengths when noise is added to both \textit{elev} and \textit{popd}.}
\label{fig_noise}
\end{figure}

On the synthetic dataset with varying levels of added noise, the GCMC results are presented in Figure~\ref{fig_noise}. As the noise ratio increases, the causal strength from elev to popd (elev \(\rightarrow\) popd) gradually decreases. Nonetheless, even with $100\%$ added noise, the causal strength of the primary causal direction remains statistically significant and still exceeds the original noise-free threshold of 0.1. In contrast, the causal strength from popd to elev (popd \(\rightarrow\) elev), representing the secondary direction, remains relatively stable but consistently below the 0.1 threshold. These results highlight the ability of GCMC to distinguish dominant from secondary causal directions under varying noise conditions. Theoretically, GCMC quantifies causal strength based on the number of shared nearest neighbors in mutual cross mapping, rather than predictive skill as used in GCCM. This mechanism provides GCMC with enhanced robustness to observational noise. Overall, the experiments on synthetic data confirm that GCMC is effective in capturing meaningful causal relationships, even in the presence of observational noise.

\section{Discussion}\label{sec4} 

The Earth system is a dynamic system, and the empirical dynamic modeling (EDM) framework, which is applicable to dynamic systems theory\citep{ma2017cme,leng2020pcm,tao2023cmc}, provides an effective tool for causal inference in Earth system science\citep{gao2023gccm}. By appropriately extending time-series causation models within the EDM framework to spatial cross-sectional data, we can effectively examine causal inference issues in Earth system science from a spatial perspective. The case studies in this paper involve four commonly used models for spatial cross-sectional data. Among them, GCCM and GCMC are spatial causation models designed based on causal principles. GCCM extends CCM to the spatial dimension, while in this study, GCMC extends CMC to the spatial dimension. The correlation measured by PCC helps us understand causation and is also the method used within GCCM to quantify causal associations. The GD model, grounded in the principle of spatial stratified heterogeneity \citep{wang2010gd1,wang2016gd2,lyu2025gdverse}, which, to some extent, implies causation through the strength of detected spatial associations \citep{wang2024gd3}. However, in our case studies, the strength of spatial association detected by GD fails to distinguish primary and secondary causal directions in certain scenarios. PCC and GD are not designed based on rigorous causal principles. While they can aid in understanding causation but cannot be directly applied to data-driven causal discovery. 

There are considerable differences between GCCM and GCMC, which are summarized in Table~\ref{tbl_diff}. Both GCCM and GCMC require measuring causal strengths within the reconstructed state space \citep{Takens1981embedding,Mane1981embedding,sugihara2012ccm}. In this process, GCCM needs preprocessing to remove linearity \citep{gao2023gccm}, whereas GCMC does not require the removal of linear trends. The decision to remove linearity or not only affects the magnitude of the causal strength measured by GCMC. This distinction arises from the fact that GCCM and GCMC employ different approaches to quantify causal strengths. GCCM uses prediction skill to measure causation \citep{sugihara2012ccm,gao2023gccm}, while GCMC relies on the intersectional cardinality of nearest neighbors. In determining specific causation, GCCM utilizes the convergence of prediction skill \citep{sugihara2012ccm,gao2023gccm}, whereas GCMC depends on the statistical significance of the estimated causal strength. For bivariate causal diagnostics, both GCCM and GCMC can effectively identify the primary causal direction, while the secondary causal direction requires further determination to ascertain whether it represents reverse causation or a reactive relationship of the primary causal direction. Similar to GCCM, GCMC is an extension of empirical dynamic modeling frameworks applied to spatial cross-sectional data, but it approaches causal inference in the spatial dimension from a different manner. Compared to GCCM, GCMC offers computational advantages. GCCM measures causation through the convergence of prediction skill, which requires performing calculations with different numbers of spatial units in varying windows across the spatial cross-section (i.e., using certain libraries). In contrast, GCMC only requires performing intersectional statistics on a fixed number of nearest neighbors. GCMC can be used in conjunction with GCCM to enhance causal inference driven by spatial cross-sectional data.

\begin{table}[htb]
\centering
\caption{A comparison of characteristics of GCCM and GCMC in measuring causal strength}
\label{tbl_diff}
{\scriptsize
\begin{tabular}{|c|c|c|}
\hline
\textbf{Model characteristics} & \textbf{GCCM} & \textbf{GCMC} \\
\hline
Removing the linearity in advance & Need & No need \\
\hline
The causation concept & Predictability & Intersectional cardinality \\
\hline
Measure of causal strength & Pearson correlation coefficient & DeLong placement \\ 
\hline
Identification of causation & Convergence of prediction skill \(\rho\) & Significance of causal strength \\
\hline
The secondary causation &
Requiring additional consideration & Requiring additional consideration \\
\hline
\end{tabular}
}
\end{table}

When using GCMC to measure causal strengths, the number of nearest neighbors \(k\) selected is an extremely important parameter. Taking two spatial cross-sectional variables, \(x\) and \(y\), as an example, when the number of neighbors chosen is small, the corresponding IC value will be small, which leads to a measured causal strength close to zero, making it impossible to effectively infer the relationship between \(x\) and \(y\). On the other hand, when the number of neighbors is too large, consider an extreme case where all neighbors are selected for the calculation. In this case, the IC value for both \(x \to y\) and \(y \to x\) will be 1, rendering the computed causal strength meaningless. For spatial cross-sectional data with different sample sizes, a suitable \(k\) must be selected to correctly apply GCMC for measuring causal strengths.

In this manuscript, we have selected three cases with clear causations to demonstrate the use of GCMC. Among these, cases of residential crime and population density study involve spatial lattice data, while the npp case uses spatial grid data, covering scenarios of weak, medium, and strong coupling. In the crime case, the causal strength measured by GCMC is relatively small due to the limited sample size (49 spatial units), which results in fewer available neighbors. However, GCMC still effectively measures the causal strengths between variables in the crime case. In the population density and npp cases, the sample size is two orders of magnitude larger than the crime case, and at this point, the statistical test for causal strength used internally by GCMC proves useful. The appropriate sample size is the second important prerequisite for applying GCMC. 

We also tested the ability of GCMC to handle noise in observational data and the results demonstrate that GCMC can still effectively measure causal associations even when different levels of noise are present in the observational data. Nevertheless, GCMC was not designed to distinguish between direct causal associations and indirect, transitive causation. In Earth system, many causal associations arise through causal transmission \citep{runge2019esstsc,runge2023tsc}. However, neither GCCM nor GCMC can currently separate direct and indirect causation. Future research can focus on this separation to develop new spatial causation models.

\section{Conclusion}\label{sec5} 

This study presents the geographical cross mapping cardinality (GCMC) model, a novel approach for measuring causal associations between variables from spatial cross-sectional data. It addresses a long-standing challenge in GIScience and Earth system science: When time series-based causal inference methods prove inadequate, models capable of inferring causality from a spatial perspective are often absent. GCMC transforms the causal associations, which are difficult to measure in the original state space, into a lagged state space where empirical dynamic modeling theory can be applied, thus enabling inference of causation in systems lacking explicit temporal variation. In the reconstructed state space, GCMC characterizes causation based on the number of intersections of nearest neighbors in mutual cross mapping, quantifies causal strength using the DeLong method, and calculates its statistical significance and confidence intervals. This approach contributes a novel, data-driven framework that extends the applicability of empirical dynamic modeling to spatial domains, enhancing the analytical toolkit available for causal inference in GIScience and Earth system science. The applicability and generalizability of GCMC are demonstrated through three case studies with varying sample sizes, data formats, and coupling conditions, while its sensitivity to noise is evaluated using controlled synthetic data. The proposed approach yields an interpretable and robust tool for causal inference in spatial cross-sectional data, facilitating deeper understanding of causal mechanisms underlying geographic phenomena. However, GCMC focuses on bivariate relationships and depends on the selection of an appropriate number of neighbors to measure causal associations, thus constraining its broader applicability. Future research can expand its scope by distinguishing between direct and indirect causal associations and developing automated strategies for identifying complex causation in spatial data.

\section*{Disclosure statement}

No conflict of interest exists in this manuscript, and the manuscript was approved by all authors for publication.

% \section*{Funding}

\section*{Data availability statement}

All data and codes supporting the findings of this study are publicly available at \\
\url{https://figshare.com/s/a1a9af2cafdf33581c55} .

% All the data used in this study are available from the spEDM package repository 
% (\url{https://github.com/stscl/spEDM}), specifically within the \texttt{inst/case} directory. 
% The code supporting the findings of this study is publicly available at 
% \url{https://github.com/SpatLyu/GCMC}.

%\section*{References}
\baselineskip12pt
\bibliographystyle{elsarticle-harv} 
\bibliography{references.bib} 



\end{document}
